\documentclass[xcolor={svgnames}]{beamer}
\usetheme[
numbering=fraction % curr_page / total_page
,block=fill        % 为 block 显示方框
% ,progressbar=frametitle % add a progressbar at frametitle
,sectionpage=progressbar
]{metropolis}

%% 保留熟悉的数学字体：两种方法
% \usefonttheme[onlymath]{serif}
\renewcommand\mathfamilydefault{\rmdefault}


\input ../defs.tex

% \usepackage[x11names]{xcolor}

\usepackage[english]{babel}
% \usepackage{enumitem} % 支持自定义 item symbol
\usepackage{tikz}
\usetikzlibrary{tikzmark,positioning,shapes,calc}
\usepackage{mathtools}
\usepackage{fancybox} %支持文本框
\usepackage{subfig} %支持subfig
% add following for otherwise it will complain
% see https://tex.stackexchange.com/questions/426088/texlive-pretest-2018-beamer-and-subfig-collide
\makeatletter
\let\@@magyar@captionfix\relax
\makeatother
%%%

%=== font customization ===%
%%
%% Note that once use customizing fonts, switch to xelatex
%%
% \usepackage{fontenc}
% \setsansfont{Varela Round}
% % \setsansfont{IM FELL DW Pica}
% \setmonofont{DejaVu Sans Mono}
% \setmathfont{Fira Sans}


% % 目录标数字
% \setbeamertemplate{section in toc}[sections numbered]
% % 无序列表用实心点
% \setbeamertemplate{itemize item}{$\bullet$}
% \setbeamertemplate{navigation symnols}{}

%% 设置列表的 item symbol
\setbeamertemplate{itemize item}{\textbullet}
\setbeamertemplate{itemize subitem}{$\blacktriangleright$}
\setbeamertemplate{itemize subsubitem}{$\checkmark$}

%% 定理环境计数
\setbeamertemplate{theorems}[numbered]
% \setbeamertemplate{lemmas}[numbered]

% 定义颜色
\definecolor{alizarin}{rgb}{0.82, 0.1, 0.26} % 红色
\definecolor{DarkFern}{HTML}{407428} % 绿色
\definecolor{tyorange}{HTML}{ff6600}
\definecolor{tyblue}{HTML}{0066ff}
\colorlet{main}{DarkFern!100!white} % 第一种设置方法
\colorlet{main}{red!70!black} % 第二种设置方法
\definecolor{bistre}{rgb}{0.24, 0.17, 0.12} % 黑色
\definecolor{mygrey}{rgb}{0.52, 0.52, 0.51} % 灰色
\colorlet{main}{tyblue}
\colorlet{text}{bistre!100!white}

% 不同元素指定不同颜色，fg是本身颜色，bg是背景颜色，!num!改变数值提供渐变色
% \setbeamercolor{title}{fg=main}
% \setbeamercolor{frametitle}{fg=main}
% \setbeamercolor{section in toc}{fg=text}
% \setbeamercolor{normal text}{fg=text}
% \setbeamercolor{block title}{fg=alizarin,bg=mygrey!20!white}
% \setbeamercolor{block body}{fg=black,bg=mygrey!15!white}
% \setbeamercolor{qed symbol}{fg=main} % 证明结束后的框颜色
% \setbeamercolor{math text}{fg=black}

%=== set theme background color to white ===%
% \setbeamercolor{background}{bg=black}
\setbeamercolor{alerted text}{fg=red!50!black}
\setbeamercolor{normal text}{fg=black!85!}
\setbeamercolor{example text}{fg=black!50!cyan}
% \setbeamercolor{frametitle}{bg=Grey}



%-------------------正文-------------------------%

\title{Quickest Search Over Multiple Sequences}
\subtitle{Group Meeting}
\author{Bingbing Hu}
\institute{SIST, ShanghaiTech}
\date{\today}

\begin{document}

\frame[plain,noframenumbering]{\titlepage}

%=== index page ===%
\begin{frame}{Outline}
\tableofcontents
\end{frame}

%=== begin your slides ==%
\section{Motivation}
%% 生成当前 section 的 outline page
% \frame{\frametitle{Outline}\tableofcontents[currentsection]}

\begin{frame}{Motivation}
  It has many applicaitons:
  \begin{itemize}
    \item Locate attack point in sensor networks
    \item Find free frequency channel in cognitive radio systems
    \item Quality monitoring of manufacturing machines
    \item Sequentially search for certain data over multiple databases
  \end{itemize}
\end{frame}

\section{Model}

% \frame{\frametitle{Outline}\tableofcontents[currentsection]}

\begin{frame}{Model}
  \framesubtitle{One more subtitle}  % disabled in metropolis
  Consider $N$ sequences

  \begin{itemize}
    \item $N$ sequences: $\{Y_k^i, ~k =1,2,\dots\}, ~i=1,\dots,N$
    \item For each $i$: $\{Y_k^i, ~k =1,2,\dots\}$ are i.i.d. observations
    \item Two hypotheses:
    $$
      \begin{array}{rl}
        1-\pi_0 &\quad H_0: ~Y_k^i \sim Q_0, ~k=1,2,\dots \\
        \tikzmark{md1}\pi_0 &\quad H_1: ~Y_k^i \sim Q_1, ~k=1,2,\dots
      \end{array}
    $$
  \end{itemize}
  \begin{tikzpicture}[
    remember picture,
    overlay,
    expl/.style={draw=orange,fill=orange!30,rounded corners},
    arrow/.style={red!80!black, thick,->,>=latex}
    ]
    \node<2>[expl]
    (md1ex)
    at (2,-0.5)
    {prior prob.};

    \draw<2>[arrow]
    (md1ex) to[out=90,in=180] ([yshift=0.5ex]{pic cs:md1});
  \end{tikzpicture}
  \onslide<3->

  \bigskip
  Intuitively, we want to quickly search the seq. generated by $Q_1$ with \alert{\textbf{minimal delay}} and \alert{\textbf{maximal accuracy}}.
\end{frame}
%%%%%%%%%%%%%%%%%%%%%%%%%%%%%%%%%%%%%%%%%%%%%%%%%%%%%%%
\begin{frame}{Model (cont.)}
  \begin{itemize}
    \item Each time select a sequence, say seq. $j$
    \item Draw a smaple $Y$ from selected seq.
    \item Make one of the actions:
    \begin{itemize}
      \item[$\blacktriangleright$] \textbf{\alert{Decision}}: stop and claim
      \item[$\blacktriangleright$] \textbf{\alert{Observation}}: continue sampling from the same seq.
      \item[$\blacktriangleright$] \textbf{\alert{Exploration}}: observe next seq.
    \end{itemize}

    \item Stop when proper conditions are satisfied
  \end{itemize}
  \pause
  % \bigskip
  % Remark: switch back is not allowed
  Remarks:
    \begin{itemize}\itemsep-2pt
      \item Consider $N = \infty$
      \item Switch back is not allowed
      \onslide<3->
      \item A1: All seq.s are independent
      \item<4-> A2: A seq.'s distribution is independent of all others'
    \end{itemize}
\end{frame}
%%%%%%%%%%%%%%%%%%%%%%%%%%%%%%%%%%%%%%%%%%%%%%%%%%%%%%%
\begin{frame}{Notations}
  \begin{itemize}
    \item Time $k = 1,2,\dots$
    \item $s_k$: seq. index observed at time $k$
    \item $Y_k^{s_k}$: the observation at time $k$
    \item Filtration $\mathcal{F}_k = \sigma(Y_1^{s_1},Y_2^{s_2},\dots,Y_k^{s_k})$: knowledge up to time $k$
    \item Stopping time $\tau$: stop at seq. $s_{\tau}$ and claim
    \item $\mathcal{T}$: the set of all stopping times w.r.t. $\mathcal{F}_k$
    \item Switching function
    $$
      \phi_k(\mathcal{F}_k) =
      \begin{cases}
        1, &\quad \text{\alert{exploration}}, s_{k+1} = s_k + 1 \\
        0, &\quad \text{\alert{observation}}, s_{k+1} = s_k
      \end{cases}
    $$
  \end{itemize}
\end{frame}
%%%%%%%%%%%%%%%%%%%%%%%%%%%%%%%%%%%%%%%%%%%%%%%%%%%%%%%%
\begin{frame}{Problem}
  Two performance metrics
  \begin{itemize}
    \item \textbf{False alarm:} $P(H^{s_{\tau}} = H_0)$
    \item \textbf{Avg. \# of samples:} $\E(\tau)$
  \end{itemize}
  This leads to the core problem
  \begin{block}{Problem}
    \begin{equation}\label{eq:origin}
      \inf_{\tau \in \mathcal{T}, \phi}\left[P(H^{s_{\tau}} = H_0) + c\E(\tau)\right]
    \end{equation}
  \end{block}
\end{frame}
%%%%%%%%%%%%%%%%%%%%%%%%%%%%%%%%%%%%%%%%%%%%%%%%%%%%%%%

\section{Solutions}
\subsection{Finite Horizon}
\begin{frame}{Finite Horizon}
  \bf
  First, we consider a finite-horizon ver. of Problem \eqref{eq:origin} by restricting \alert{$\tau \in [0, T]$}. This means we must stop at time $T$.
\end{frame}
%%%%%%%%%%%%%%%%%%%%%%%%%%%%%%%%%%%%%%%%%%%%%%%%%%%%%%%%%%%%%%%%
\begin{frame}{Posterior Prob.}
  \onslide<1>{
  Define $\pi_k = P(H^{s_k} = H_1 ~|~ \mathcal{F}_k)$ as the posterior prob. that seq. $s_k$ is generated by $Q_1$ after observing $\mathcal{F}_k$,}
  \begin{equation}\label{eq:core-recur}
    \begin{split}
      \pi_1 =& \frac{\tikzmark{pi0}\pi_0 q_1(Y_1^1)\tikzmark{q1}}{\pi_0 q_1(Y_1^1) + (1-\pi_0) q_0(Y_1^1)} \\
      &\vdots \\
      \pi_{k+1} =& \frac{\pi_k q_1(Y_{k+1}^{s_{k+1}})}{\pi_k q_1(Y_{k+1}^{s_{k+1}}) + (1-\pi_k) q_0(Y_{k+1}^{s_{k+1}})} \ones(\phi_k=0) \\
      &+ \frac{\pi_0 q_1(Y_{k+1}^{s_{k+1}})}{\pi_0 q_1(Y_{k+1}^{s_{k+1}}) + (1-\pi_0) q_0(Y_{k+1}^{s_{k+1}})} \ones(\phi_k=1)
    \end{split}
  \end{equation}
  \begin{tikzpicture}[
    remember picture,
    overlay,
    expl/.style={draw=orange,fill=orange!30,rounded corners},
    arrow/.style={red!80!black, thick,->,>=latex}
    ]
    \node<2->[expl]
    (pi0ex)
    at (2.5,5.5)
    {$P(H^1 = H_1)$};
    \node<3->[expl]
    (q1ex)
    at (6,5.5)
    {$P(Y^1_1 ~|~ H^1 = H_1)$};

    \draw<2->[arrow]
    (pi0ex) to ([yshift=0.5ex]{pic cs:pi0});
    \draw<3->[arrow]
    (q1ex) to  ([yshift=0.5ex]{pic cs:q1});

  \end{tikzpicture}
\end{frame}
%%%%%%%%%%%%%%%%%%%%%%%%%%%%%%%%%%%%%%%%%%%%%%%%%%%%%
\begin{frame}{Cost-to-go Function}
  \textbf{At each time $k~(\le T)$, based on $\mathcal{F}_k$, shall we stop sampling?}

  \bigskip
  \only<2>{This calls the concept of function $\tilde{J}^T_{k}(\mathcal{F}_k)$, standing for the minimal cost-to-go at time $k$.}

  \only<3->{More precisely, \alert{at time $k$, based on $\mathcal{F}_k$, take action (stop or continue?), pay a cost $\tilde{J}^T_k$}.}

  \visible<2->{
  \begin{equation}\label{eq:cost-to-go}
    \displaystyle \tilde{J}^T_{k}(\mathcal{F}_k) = \min\left\{
    \underbrace{1 - \pi_k}_{\text{\alert{stop cost}}},
    ~\underbrace{c + \inf_{\phi_k}\E\left\{
                    \tilde{J}^T_{k+1}(\mathcal{F}_{k+1}) ~|~ \mathcal{F}_k, \phi_k
                  \right\}}_{\text{\alert{cont. cost}}}
    \right\}
  \end{equation}
  }
  % \visible<2>
  \begin{itemize}
    \item<4-> e.g., $\tilde{J}^T_{T}(\mathcal{F}_T) = 1 - \pi_T$
  \end{itemize}
\end{frame}
%%%%%%%%%%%%%%%%%%%%%%%%%%%%%%%%%%%%%%%%%%%%%%%%%%
\begin{frame}{Finite Horizon}
  \begin{lemma}
    For each $k$, function $\tilde{J}^T_{k}(\mathcal{F}_k)$ depends only on $\pi_k$, say $J^T_{k}(\pi_k)$, so as the optimal switching rules $\phi_k$.
  \end{lemma}
  \pause

  To prove the lemma
  \begin{itemize}[<+->]
    \item[\checkmark] $\tilde{J}^T_{T}(\mathcal{F}_T) = 1 - \pi_T$
    \item Assume $\tilde{J}^T_{k+1}(\mathcal{F}_{k+1})$ depends on $\pi_{k+1}$ $\implies$ $\tilde{J}^T_{k}(\mathcal{F}_{k})$ depends only on $\pi_k$
    \item So as $\phi_k$
  \end{itemize}
\end{frame}
%%%%%%%%%%%%%%%%%%%%%%%%%%%%%%%%%%%%%%%%%%%%%%%%%
\begin{frame}{Proof of Lemma 1}
  Recall (\ref{eq:cost-to-go}),
  $$
  \tilde{J}^T_{k}(\mathcal{F}_k) = \min\left\{
  1 - \pi_k,
  ~c + \inf_{\phi_k}\E\left\{
                  \tilde{J}^T_{k+1}(\mathcal{F}_{k+1}) ~|~ \mathcal{F}_k, \phi_k
                \right\}
  \right\}.
  $$
  The second item can be rewritten
  \[
    c + \min\{A^T_{k,c}(\pi_k), A^T_{k,s}\},
  \]
  so $\tilde{J}^T_{k}(\mathcal{F}_k)$ depends on $\pi_k$, and $\phi_k$ is decided by
  \begin{equation}
    \phi_k =
    \begin{cases}
      1, &\quad \text{if } A^T_{k,c}(\pi_k) > A^T_{k,s} \\
      0, &\quad \text{otherwise}
    \end{cases},
  \end{equation}
  which also depends only on $\pi_k$.
\end{frame}
%%%%%%%%%%%%%%%%%%%%%%%%%%%%%%%%%%%%%%%%%%%%%%%%%
\begin{frame}{Proof of Lemma 1 (cont.)}
\[
  \begin{gathered}
  \tilde{J}^T_{k}(\mathcal{F}_k) = \min\left\{
  1 - \pi_k,
  ~c + \alert{\Ovalbox{$\inf_{\phi_k}
    \E\left\{
      \tilde{J}^T_{k+1}(\mathcal{F}_{k+1}) ~|~ \mathcal{F}_k, \phi_k
    \right\}$}}
  \right\} \\[3pt]
  \implies
  \min\left\{
    \E_{\phi_k=0}(\dots), ~\E_{\phi_k=1}(\dots)
  \right\} \\
  \implies
  \min\left\{
    \int J^T_{k+1}(\pi_{k+1}) \tikzmark{fc} f_c(y|\mathcal{F}_k) dy,~
    \int J^T_{k+1}(\pi_{k+1}) \tikzmark{fs} f_s(y|\mathcal{F}_k) dy
  \right\}
\end{gathered}
\]
\onslide<2>{
Two conditional PDF:
\begin{itemize}
  \item \textbf{stay:} $f_c(y|\mathcal{F}_k) = \pi_k q_1(y) + (1-\pi_k)q_0(y)$
  \item \textbf{switch:} $f_s(y|\mathcal{F}_k) = \pi_0 q_1(y) + (1-\pi_0)q_0(y)$
\end{itemize}
Additionally, $\pi_{k+1}$ depends on $\pi_k$ or $\pi_0$ by \eqref{eq:core-recur}.
}
\begin{tikzpicture}[
  remember picture,
  overlay,
  expl/.style={draw=orange,fill=orange!30,rounded corners},
  arrow/.style={red!80!black, thick,->,>=latex}
  ]
  \node<3->[expl]
  (fcex)
  at (-3.45,2)
  {$A^T_{k,c}(\pi_k)$};
  \node<3->[expl]
  (fsex)
  at (1,2)
  {$A^T_{k,s}$};

  \draw<3->[arrow]
  (fcex) to ([yshift=-1ex]{pic cs:fc});
  \draw<3->[arrow]
  (fsex) to  ([yshift=-1ex]{pic cs:fs});
\end{tikzpicture}
% \only<4->{
% % \bigskip
% \begin{itemize}
%   \item $A^T_{k,c}(\pi_k)$: expected next step cost when stay
%   \item $A^T_{k,s}$: expected next step cost when switch
% \end{itemize}
% }
\end{frame}
%%%%%%%%%%%%%%%%%%%%%%%%%%%%%%%%%%%%%%%%%%%%%%%%%%%%%%%
\begin{frame}{Finite Horizon}
  \begin{lemma}\label{lem2}
    The functions $J^T_{k}(\pi_k)$ and $A^T_{k,c}(\pi_k)$ are nonnegative concave w.r.t. $\pi_k \in [0,1]$. And $J^T_{k}(1) = A^T_{k,c}(1) = 0$.
  \end{lemma}
  \begin{itemize}
    \item These two functions are both somewhat nonnegative costs
    \item Concavity reserved
    \item Since $J^T_{k}(1) = \min\{ 0, \text{something}\} \ge 0 \implies J^T_k(1) = 0$, and
    \[
      A_{k,c}^T(1) = \E\left\{J^T_{k+1}(1) ~|~\mathcal{F}_k, \phi_k=0\right\} = 0
    \]
  \end{itemize}
\end{frame}
%%%%%%%%%%%%%%%%%%%%%%%%%%%%%%%%%%%%%%%%%%%%%%%%%%%%%%%%
\begin{frame}{Finite Horizon}
  \begin{columns}
    \column{.4\textwidth}
    $J^T_k(\pi_k)$ satisfy:
    \bigskip
    \begin{itemize}
      \item $\le 1-\pi_k$
      \item $J^T_k(1) = 0$
      \item concave w.r.t. $\pi_k$
    \end{itemize}
    \pause

    \column{.6\textwidth}
    \begin{figure}
        \centering
        \includegraphics[width=1.0\textwidth]{fig1}\\[-8pt]
        {\tiny src: L Lai, et al. "Quickest search"}
        \caption{An illustration of $J^T_k(\pi_k)$}
        \label{fig1}
    \end{figure}
  \end{columns}
\end{frame}
%%%%%%%%%%%%%%%%%%%%%%%%%%%%%%%%%%%%%%%%%%%%%%%%%%%
\begin{frame}{Finite Horizon (solution)}
  \begin{theorem}
    For the finite-horizon version of Problem (\ref{eq:origin}), the optimal stopping time is \alert{$\tau_{\text{opt}} = \inf\{k: \pi_k > a_k^T\}$}, in which $a_k^T$ is given by the following equation:
    \[
      1 - a_k^T = c + \min\{A_{k,c}^T(a_k^T), ~A_{k,s}^T\}.
    \]
  \end{theorem}
  \only<2>{
  Recall that
  \[
  J^T_{k}(\pi_k) = \min\left\{
  1 - \pi_k,
  ~c + \min\{A^T_{k,c}(\pi_k), A^T_{k,s}\}
  \right\}.
  \]
  }

  \only<3>{
  Once $\pi_k > a^T_k$, we have
  \[
    1 - \pi_k < 1-a^T_k = c + \min\{..\}
  \]
  which means \alert{stop cost} $<$ \alert{cont. const}. $\implies$ \textbf{stop!}
  }
\end{frame}
%%%%%%%%%%%%%%%%%%%%%%%%%%%%%%%%%%%%%%%%%%%%%%%%%%
\subsection{Infinite Horizon}

\begin{frame}{Infinite Horizon}
  \textbf{Now let \alert{$T \to \infty$}, we get the infinite horizon ver. of Problem \eqref{eq:origin}.}
\end{frame}
%%%%%%%%%%%%%%%%%%%%%%%%%%%%%%%%%%%%%%%%%%%%%%%%%
\begin{frame}{Infinite Horizon}
  Observing that
  \begin{itemize}
    \item \textbf{Constraint relaxed}: $J^{T+1}_k(\pi) \le J^{T}_k(\pi)$
    \item \textbf{Lower bounded}: $\displaystyle \lim_{T\to \infty} J^T_k(\pi) = \inf_{T>k} J^T_k(\pi) = J^{\infty}_k(\pi)$
    \item \textbf{Concavity}: $J^T_k(\pi)$ is concave for each $T, ~k$
    \item \textbf{I.I.D. observation}: $J^{\infty}_k(\pi) = J^{\infty}_{k+1}(\pi)$?
  \end{itemize}
  \pause
  % we now \eqref{eq:test} and normal \ref{eq:test}
  $\implies J(\pi):= J^{\infty}_{k}(\pi)$ concave w.r.t. $\pi$
\end{frame}
%%%%%%%%%%%%%%%%%%%%%%%%%%%%%%%%%%%%%%%%%%%%%%%%
\begin{frame}{Infinite Horizon}
  \begin{figure}
      \centering
      \includegraphics[width=.6\textwidth]{fig4}\\[-8pt]
      {\tiny src: L Lai, et al. "Quickest search"}
      \caption{An illustration of $J(\pi)$}
  \end{figure}
\end{frame}
%%%%%%%%%%%%%%%%%%%%%%%%%%%%%%%%%%%%%%%%%%%%%%%%%%%%%
\begin{frame}{Infinite Horizon}
  As $T\to \infty$,
  \begin{itemize}
    \item $A^T_{k,c}(\pi) \to A_c(\pi)$
    \item $A^T_{k,s} \to A_s$
  \end{itemize}

  \begin{lemma}\label{acas}
    \begin{equation}
      A_c(\pi)
      \begin{cases}
        > A_s, \quad \text{if } \pi < \pi_0 \\
        = A_s, \quad \text{if } \pi = \pi_0 \\
        < A_s, \quad \text{if } \pi > \pi_0
      \end{cases}
    \end{equation}
  \end{lemma}
\end{frame}
%%%%%%%%%%%%%%%%%%%%%%%%%%%%%%%%%%%%%%%%%%%%%%%%%%%%%
\begin{frame}{Infinite Horizon}
  \begin{columns}
    \column{.4\textwidth}
    \begin{itemize}
      \item $A_c(0) > A_s$
      \item $A_c(1) = 0 < A_s$
      \item $A_c(\pi)$ concave
    \end{itemize}
    \pause

    \column{.6\textwidth}
    \begin{figure}
        \centering
        \includegraphics[width=1.0\textwidth]{fig2}\\[-8pt]
        {\tiny src: L Lai, et al. "Quickest search"}
        \caption{An illustration of $A_c(\pi)$ and $A_s$}
        \label{fig2}
    \end{figure}
  \end{columns}
\end{frame}
%%%%%%%%%%%%%%%%%%%%%%%%%%%%%%%%%%%%%%%%%%%%%%%%%%%%%%%%
\begin{frame}{Infinite Horizon (solution)}
  \begin{theorem}
    The optimal stopping time for Problem (\ref{eq:origin}) is given by \alert{$\tau_{\text{opt}} = \inf\{k: \pi_k > \pi_U^*\}$} in which
    \[
      1 - \pi_U^* = c + \min\{A_{c}(\pi_U^*), ~A_{s}\}.
    \]
  \end{theorem}
  \only<2>{
  Recall that
  \begin{equation}\label{eq:core-inf}
    J(\pi) = \min\left\{
    1 - \pi,
    ~c + \min\{A_c(\pi), A_s\}
    \right\}.
  \end{equation}
  }

  \only<3>{
  Once $\pi_k > \pi_U^*$, we have
  \[
    1 - \pi_k < 1 - \pi_U^* = c + \min\{..\}
  \]
  which means \alert{stop cost} $<$ \alert{cont. const}. $\implies$ \textbf{stop!}
  }
\end{frame}
%%%%%%%%%%%%%%%%%%%%%%%%%%%%%%%%%%%%%%%%%%%%%%%%%%%%%
\begin{frame}{Process Tracking}
  \begin{figure}
      \centering
      \includegraphics[width=.85\textwidth]{fig3}\\[-8pt]
      {\tiny src: L Lai, et al. "Quickest search"}
  \end{figure}
\end{frame}
%%%%%%%%%%%%%%%%%%%%%%%%%%%%%%%%%%%%%%%%%%%%%%%%%%
\begin{frame}{For Convenience}
  \begin{figure}
  \centering
  \subfloat{
  \includegraphics[width=0.5\textwidth]{fig4}
  }
  % \quad
  \subfloat{
  \includegraphics[width=0.5\textwidth]{fig2}
  }
  \end{figure}
\end{frame}

% \section{Sample frame title}
%
% % \frame{\frametitle{Outline}\tableofcontents[currentsection]}
%
% \begin{frame}
% \begin{exampleblock}{Baye's theorem}
% \[
% \tikzmark{ptd}p(\theta\, |\, D) = \frac{\tikzmark{pdt}p(D\,|\,\theta) p(\theta)}{\tikzmark{pd}p(D)}
% \]
% \begin{tikzpicture}[
%   remember picture,
%   overlay,
%   expl/.style={draw=orange,fill=orange!30,rounded corners,text width=3cm},
%   arrow/.style={red!80!black,ultra thick,->,>=latex}
% ]
% \node<2->[expl]
%   (ptdex)
%   at (2,-2cm)
%   {Some explanation};
% \node<3->[expl]
%   (pdtex)
%   at (6,3.5cm)
%   {Some other explanation; this one is a little longer};
% \node<4->[expl]
%   (pdex)
%   at (9,-3cm)
%   {Some other explanation};
% \draw<2->[arrow]
%   (ptdex) to[out=100,in=180] ([yshift=0.5ex]{pic cs:ptd});
% \draw<3->[arrow]
%   (pdtex.west) to[out=180,in=180] ([yshift=0.5ex]{pic cs:pdt});
% \draw<4->[arrow]
%   (pdex.north) to[out=90,in=180] ([yshift=0.5ex]{pic cs:pd});
% \end{tikzpicture}
% \end{exampleblock}
% \end{frame}
%
%
% \begin{frame}
% \frametitle{Sample frame title}
% This is a text in second frame.
% For the sake of showing an example.
%
% \begin{itemize}
%  \item<1-> Text visible on slide 1
%  \item<2> Text visible on slide 2
%  \item Text visible on slide 3
% \end{itemize}
%
% % \vspace{0.3cm}
% \bigskip
%
% another example
%
% \begin{itemize}\itemsep2em
%  \item Text visible on slide 1
%  \item Text visible on slide 2
%  \item Text visible on slide 3
% \end{itemize}
% \end{frame}
%
% \begin{frame}
%   \begin{proof}
%     $$
%  \frac{1}{\displaystyle 1+
%    \frac{1}{\displaystyle 2+
%    \frac{1}{\displaystyle 3+x}}} +
%  \frac{1}{1+\frac{1}{2+\frac{1}{3+x}}}
% $$
% $$\int_0^\infty e^{-x^2} dx=\frac{\sqrt{\pi}}{2}$$
% \begin{equation} x=y+3 \label{eq:xdef}
% \end{equation}
% In equation (\ref{eq:xdef}) we saw $\dots$
%   \end{proof}
% \end{frame}
%
% \begin{frame}[label=simmonshall]{Text blocks}
%   \framesubtitle{In plain, example, and \alert{alert} flavour}
%   \alert{This text} is highlighted.
%
%   \begin{block}{A plain block}
%     This is a plain block containing some \alert{highlighted text}.
%   \end{block}
%   \begin{exampleblock}{An example block}
%     This is an example block containing some \alert{highlighted text}, and this is \emph{emphasized}.
%   \end{exampleblock}
%   \begin{alertblock}{An alert block}
%     This is an alert block containing some \alert{highlighted text}.
%   \end{alertblock}
% \end{frame}

\begin{frame}[standout]
  \Huge Thank you!
\end{frame}

% \begin{frame}[label=conclusion, standout]{Conclusion}
%   Awesome slide
% \end{frame}


\end{document}
